\documentclass[STT-4300-Main.tex]{subfiles}
\title{Chapitre 1 - Introduction et motivation}
\subtitle{STT-4300 - Régression II}
\date{\today}
\begin{document}
\begin{refsection}

\begin{frame}[t]
\titlepage
\end{frame}

\section{Introduction et motivation}

\begin{frame}{}
\begin{center}
{\color{blue}«~Statistics is often wrongly perceived as a set of tools instead of a set of problems~»}
\end{center}

\bigskip
\bigskip

tiré de J.~H.~Friedman (2001) ``The Role of Statistics in Data Revolution'' \textit{International Statistical Review}, 69, 5--10.
\end{frame}

\subsection{Exemple 1 : Boston Housing}

\begin{frame}{Exemple 1 -- Boston Housing}

Les données \textit{Boston Housing} ont été utilisées pour décrire la relation entre la valeur des logements dans les banlieues de Boston et différents attributs.

\bigskip

La taille de l'échantillon est de 506.

\bigskip

Les données proviennent à l'origine de \citeN{Harr:Rubi:1978}, utilisées dans
\citeN{Bels:Kuh:Wels:1980}, disponibles depuis l'archive \href{https://lib.stat.cmu.edu/datasets/boston}{StatLib}.


\end{frame}

\begin{frame}{}

\begin{table}[htbp]
  \begin{center}
{\small \begin{tabular}{ll}
\hline
\texttt{medv} & valeur médiane des maisons occupées par le propriétaire en \$1000\\
\texttt{crim} & taux de criminalité par habitant par ville\\
\texttt{zn} & proportion de terrains résidentiels zonés pour lots \\
 & de plus de 25\,000 pi$^2$\\
\texttt{indus} & proportion d'acres commerciales non liées au commerce de détail par ville\\
\texttt{nox} & concentration d'oxydes nitriques (parties par 10 millions)\\
\texttt{rm} & nombre moyen de pièces par logement\\
\texttt{age} & proportion de logements occupés par le propriétaire construits avant 1940\\
\texttt{dis} & distances pondérées vers cinq centres d'emploi de Boston\\
\texttt{rad} & indice d'accessibilité aux autoroutes radiales\\
\texttt{tax} & taux de taxe foncière à pleine valeur par \$10\,000\\
\texttt{ptratio} & ratio élèves-enseignant par ville\\
\texttt{b} & $1000(\mbox{Bk}-0.63)^2$, où Bk est la proportion \\
 & de personnes noires par ville\\
\texttt{lstat} & proportion de la population de statut inférieur\\
\texttt{chas} & variable indicatrice de la rivière Charles \\
& (=1 si le secteur borde la rivière ; 0 sinon)\\
\hline
\end{tabular}}
  \end{center}
\caption{Description des données Boston Housing.}
\label{housing.table}
\end{table}
\end{frame}


\begin{frame}[fragile]{}

Modèle~\eqref{housing.model.linear} :
\begin{eqnarray}
\label{housing.model.linear}
\lefteqn{\log(\texttt{medv}) = \beta_0 + \beta_1\texttt{crim} + \beta_2\texttt{zn}+
\beta_3\texttt{indus}+ } \nonumber \\
& + & \beta_4\texttt{nox} + \beta_5\texttt{rm} + \beta_6\texttt{age} +
\beta_7\texttt{dis} + \beta_8\texttt{rad} \nonumber \\
& + & \beta_9\texttt{tax} + \beta_{10}\texttt{ptratio} + \beta_{11}\texttt{b} +
\beta_{12}\texttt{lstat} \nonumber \\
& + & \beta_{13} \texttt{chas} + \epsilon.
\end{eqnarray}

\bigskip

{\scriptsize
\begin{verbatim}
> linearfit <- lm(log(medv)~chas+crim+zn+indus+nox+rm+age+dis+
    rad+tax+ptratio+b+lstat,data=BostonHousing)
> summary(linearfit)
\end{verbatim}
}

\end{frame}


\begin{frame}[fragile]{}

{\scriptsize
\begin{verbatim}
Coefficients:
              Estimate Std. Error t value Pr(>|t|)    
(Intercept)  4.1020423  0.2042726  20.081  < 2e-16 ***
chas         0.1008876  0.0344859   2.925 0.003598 ** 
crim        -0.0102715  0.0013155  -7.808 3.52e-14 ***
zn           0.0011725  0.0005495   2.134 0.033349 *  
indus        0.0024668  0.0024614   1.002 0.316755    
nox         -0.7783993  0.1528902  -5.091 5.07e-07 ***
rm           0.0908331  0.0167280   5.430 8.87e-08 ***
age          0.0002106  0.0005287   0.398 0.690567    
dis         -0.0490873  0.0079834  -6.149 1.62e-09 ***
rad          0.0142673  0.0026556   5.373 1.20e-07 ***
tax         -0.0006258  0.0001505  -4.157 3.80e-05 ***
ptratio     -0.0382715  0.0052365  -7.309 1.10e-12 ***
b            0.0004136  0.0001075   3.847 0.000135 ***
lstat       -0.0290355  0.0020299 -14.304  < 2e-16 ***
---
Signif. codes:  0 '***' 0.001 '**' 0.01 '*' 0.05 '.' 0.1 ' ' 1

Residual standard error: 0.1899 on 492 degrees of freedom
Multiple R-squared:  0.7896,	Adjusted R-squared:  0.7841 
F-statistic: 142.1 on 13 and 492 DF,  p-value: < 2.2e-16
\end{verbatim}
}

\end{frame}


\begin{frame}{}
\begin{center}
\includegraphics[width = 7.5cm]{figures/BHlindiag.pdf}
\end{center}
\end{frame}


\begin{frame}[fragile]{}

Modèle~\eqref{housing.model.2} (inspiré de \citeNP{Bels:Kuh:Wels:1980}) :
\begin{eqnarray}
\label{housing.model.2}
\lefteqn{\log(\texttt{medv}) = \beta_0 + \beta_1\texttt{crim} + \beta_2\texttt{zn}+
\beta_3\texttt{indus}+ }\nonumber \\
& + & \beta_4\texttt{nox}{\color{blue}^2} + \beta_5\texttt{rm}{\color{blue}^2} + \beta_6\texttt{age} +
\beta_7 {\color{blue}\log}(\texttt{dis}) + \beta_8 {\color{blue}\log}(\texttt{rad}) \nonumber  \\
& + & \beta_9\texttt{tax} + \beta_{10}\texttt{ptratio} + \beta_{11}\texttt{b} +
\beta_{12} {\color{blue}\log}(\texttt{lstat})\nonumber  \\
& + & \beta_{13} \texttt{chas} + \epsilon.
\end{eqnarray}

\bigskip

{\scriptsize
\begin{verbatim}
> linearfit2 <- lm(log(medv)~chas+crim+zn+indus+I(nox^2)+
   I(rm^2)+age+log(dis)+log(rad)+tax+ptratio+b+log(lstat),
   data=BostonHousing)
> summary(linearfit2)
\end{verbatim}
}
\end{frame}

\begin{frame}[fragile]{}

{\scriptsize
\begin{verbatim}
Coefficients:
              Estimate Std. Error t value Pr(>|t|)    
(Intercept)  4.558e+00  1.544e-01  29.512  < 2e-16 ***
chas         9.140e-02  3.320e-02   2.753 0.006129 ** 
crim        -1.186e-02  1.245e-03  -9.532  < 2e-16 ***
zn           8.016e-05  5.056e-04   0.159 0.874105    
indus        2.395e-04  2.364e-03   0.101 0.919318    
I(nox^2)    -6.380e-01  1.131e-01  -5.639 2.88e-08 ***
I(rm^2)      6.328e-03  1.312e-03   4.823 1.89e-06 ***
age          9.074e-05  5.263e-04   0.172 0.863179    
log(dis)    -1.913e-01  3.339e-02  -5.727 1.78e-08 ***
log(rad)     9.571e-02  1.913e-02   5.002 7.91e-07 ***
tax         -4.203e-04  1.227e-04  -3.426 0.000664 ***
ptratio     -3.112e-02  5.013e-03  -6.208 1.14e-09 ***
b            3.637e-04  1.031e-04   3.527 0.000460 ***
log(lstat)  -3.712e-01  2.501e-02 -14.841  < 2e-16 ***
---
Signif. codes:  0 '***' 0.001 '**' 0.01 '*' 0.05 '.' 0.1 ' ' 1

Residual standard error: 0.1825 on 492 degrees of freedom
Multiple R-squared:  0.8059,	Adjusted R-squared:  0.8008 
F-statistic: 157.1 on 13 and 492 DF,  p-value: < 2.2e-16
\end{verbatim}
}
\end{frame}


\begin{frame}{}
\begin{center}
\includegraphics[width = 7.5cm]{figures/BHlin2.pdf}
\end{center}
\end{frame}


\begin{frame}[fragile]{}

Modèle~\eqref{housing.model.np}, non paramétrique :
\begin{eqnarray*}
\label{housing.model.np}
\lefteqn{\log(\texttt{medv}) = \alpha + f_1(\texttt{crim}) + f_2(\texttt{zn})+
f_3(\texttt{indus}) +} \\
& + & f_4(\texttt{nox}) + f_5(\texttt{rm}) +  f_6(\texttt{age}) +
f_7(\texttt{dis}) + \\
& + & f_8(\texttt{rad}) + f_9(\texttt{tax}) + f_{10}(\texttt{ptratio}) \nonumber\\
& + & f_{11}(\texttt{b}) + f_{12}(\texttt{lstat}) + \beta \texttt{chas} +
\epsilon.
\end{eqnarray*}

\bigskip

{\scriptsize
\begin{verbatim}
> library(mgcv)
> BHnpfit.new <- gam(log(medv)~chas+s(crim)+s(zn)+s(indus)+ 
  s(nox,k=20)+ s(rm)+s(age)+s(dis,k=20)+s(rad,k=8)+s(tax)+s(ptratio)+
   s(b)+s(lstat), data=BostonHousing)
> plot(BHnpfit.new, scale=0, page=1)
\end{verbatim}
}

\bigskip
\bigskip

{\color{red} $\bigstar$ Sélection de variables par garrote non négatif (nonnegative garrote), voir
\citeN{Brei:1995} et
\citeN{Cant:Flem:Ronc:2011}.}
\end{frame}



\begin{frame}{}
\begin{center}
\includegraphics[width = 7.8cm]{figures/BHnpfit.pdf}
\end{center}
\end{frame}

\subsection{Exemple 2 : Diagnostic du cancer du sein}

\begin{frame}{Exemple 2 : Diagnostic du cancer du sein}

\label{example2}

Diagnostic du cancer du sein par cytologie mammaire (aspiration à l'aiguille fine plutôt que chirurgie invasive).

\bigskip
Données Breast cancer Wisconsin, disponibles depuis le University of
California at Irvine Repository Of Machine Learning Database
\url{www.ics.uci.edu/~mlearn/MLRepository.html}.

\bigskip

Variable réponse : 2 classes : bénigne (codée 0) ou maligne (codée 1).

\bigskip

Il y a 699 patientes dans l'étude, dont 16 ont une ou plusieurs valeur(s) manquante(s) dans les prédicteurs.
\end{frame}


\begin{frame}{}

Prédicteurs disponibles :
\begin{itemize}
\item  Épaisseur de l'amas (Clump Thickness) (1--10)
\item  Uniformité de la taille des cellules (Uniformity of Cell Size) (1--10)
\item  Uniformité de la forme des cellules (Uniformity of Cell Shape) (1--10)
\item  Adhésion marginale (Marginal Adhesion) (1--10)
\item  Taille des cellules épithéliales simples (Single Epithelial Cell Size) (1--10)
\item  Noyaux nus (Bare Nuclei) (1--10)
\item  Chromatine fade (Bland Chromatin) (1--10)
\item  Nucléoles normaux (Normal Nucleoli) (1--10)
\item  Mitoses (1--10)
\end{itemize}

\end{frame}



\begin{frame}[fragile]{}

On considère un modèle logistique linéaire.

\bigskip

{\scriptsize
\begin{verbatim}
> breast.fit <- glm(status~clump+cellsize+cellshape+adhesion+
singlesize+nuclei+chromatin+nucleoli+mitoses,data=breast,
family=binomial,na.action=na.omit)
> fitted.breast <- predict(breast.fit,type="response")
> plot(fitted.breast,na.omit(breast)$status,
  xlab="Fitted values",ylab="Status")
\end{verbatim}
}

\begin{center}
\includegraphics[width = 4.8cm]{figures/breastfitted.pdf}
\end{center}

\end{frame}


\subsection{Exemple 3 : Étude sur la pollution de l'air}

\begin{frame}{Exemple 3 : Étude sur la pollution de l'air}

Données provenant de la base National Morbidity, Mortality and Air Pollution Study (NMMAPS) disponible à \url{http://www.ihapss.jhsph.edu} pour
les 88 plus grandes régions métropolitaines des États-Unis.

\bigskip

La variable réponse est le nombre de décès par jour (\texttt{death}), que l'on
modélise comme une variable de Poisson.
\end{frame}

\begin{frame}{}

Variables explicatives potentiellement intéressantes :
\begin{itemize}
\item matière particulaire de moins de 10 $\mu$m de diamètre aérodynamique,
moyenne tronquée à 10\% sur tous les moniteurs d'un comté (\texttt{pm10tmean})
\item température moyenne (sur 24 heures) (\texttt{tmean})
\item temps (\texttt{day})
\item température du point de rosée (\texttt{dptp})
\item jour de la semaine (\texttt{dow} : 1=sam, 2=dim, ...)
\end{itemize}

\bigskip

L'intérêt principal porte sur l'effet de \texttt{pm10tmean} (pollution), qui entre
linéairement dans le modèle.
\end{frame}


\begin{frame}[fragile]{San Francisco, 1\textsuperscript{er} janv.\ 1987 au 31 déc.\ 2000,
$n=15342$}

{\scriptsize
\begin{verbatim}
sanFfit <- gam(death~pm10tmean+s(tmean)+s(day)+s(dptp)+factor(dow),
data=sanF,family=poisson)

 Parametric coefficients:
               Estimate Std. Error z value Pr(>|z|)    
(Intercept)   1.7820648  0.0245328  72.640  < 2e-16 ***
pm10tmean     0.0017404  0.0005446   3.196  0.00139 ** 
factor(dow)2 -0.0261762  0.0349350  -0.749  0.45369    
factor(dow)3  0.0194925  0.0343809   0.567  0.57074    
factor(dow)4 -0.0207042  0.0348386  -0.594  0.55232    
factor(dow)5 -0.0013170  0.0346407  -0.038  0.96967    
factor(dow)6 -0.0002978  0.0346081  -0.009  0.99314    
factor(dow)7 -0.0140022  0.0345012  -0.406  0.68486    
---
Signif. codes:  0 '***' 0.001 '**' 0.01 '*' 0.05 '.' 0.1 ' ' 1
\end{verbatim}
}
\end{frame}

\begin{frame}[fragile]{}

\bigskip
{\scriptsize
\begin{verbatim}
Approximate significance of smooth terms:
           edf Ref.df Chi.sq  p-value    
s(tmean) 1.973  2.568  7.642   0.0388 *  
s(day)   2.241  2.796 24.476 2.14e-05 ***
s(dptp)  1.003  1.006  0.494   0.4830    
---
Signif. codes:  0 '***' 0.001 '**' 0.01 '*' 0.05 '.' 0.1 ' ' 1

R-sq.(adj) =  0.0145   Deviance explained = 2.03%
UBRE = 0.74149  Scale est. = 1         n = 1986
\end{verbatim}
}
\end{frame}

\begin{frame}{}

\begin{center}
\includegraphics[width = 7.5cm]{figures/NMMAPSfit.pdf}
\end{center}

\end{frame}


\begin{frame}{}
\begin{center}
\includegraphics[width = 4cm]{figures/NMMAPSfitwithres.pdf}
\end{center}

Questions ouvertes :
\begin{itemize}
\item Tous les termes sont-ils significatifs?
\item Autres facteurs influençant le nombre de décès.

($R^2$ et déviance expliquée faibles car peu de variables explicatives.)
\item {\color{red} $\bigstar$ Modèles différents pour différents groupes d'âge (voir
\citeNP{Domi:Same:Zege:2000})}.
\end{itemize}
\end{frame}

\subsection{Exemple 4 : Coût de séjour pour problèmes de dos}
\label{CHUV.Ex}

\begin{frame}{Exemple 4 : Coût de séjour pour problèmes de dos}

100 patients hospitalisés au \textit{Centre Hospitalier
Universitaire Vaudois (CHUV)} à Lausanne (Suisse) durant 1999 pour
«~problèmes médicaux de dos~» (APDRG 243).

\smallskip
Données fournies par A.\ Marazzi, voir
\citeN{Mara:Yoha:2004}.

\bigskip
\bigskip

La variable réponse est le coût du séjour (\texttt{CouTot} en francs
suisses).
\end{frame}


\begin{frame}{}
Variables explicatives :
\begin{itemize}
\item durée de séjour (\texttt{LOS}, en jours)
\item type d'admission (\texttt{Typadm} : 0=planifiée, 1=urgence)
\item type d'assurance (\texttt{TypAss} : 0=régulière, 1=privée)
\item âge en années (\texttt{age})
\item sexe (\texttt{Sexe} : 0=femme, 1=homme)
\item destination à la sortie (\texttt{dest} :
1=domicile, 0=autre établissement de santé)
\end{itemize}

\bigskip

Objectifs :
\begin{itemize}
\item Identifier les facteurs influençant le coût du séjour
\item Prédiction
\end{itemize}
\end{frame}


\begin{frame}{}

\begin{center}
\includegraphics[width = 8cm]{figures/CostCHUV.pdf}
\end{center}

\end{frame}

\begin{frame}[fragile]{Modèle gamma avec lien logarithmique.}

{\scriptsize
\begin{verbatim}
Coefficients:
              Estimate Std. Error t value Pr(>|t|)    
(Intercept)  7.2338121  0.1468364  49.264  < 2e-16 ***
log(LOS)     0.8222203  0.0279641  29.403  < 2e-16 ***
Typadm       0.2136148  0.0500028   4.272 4.67e-05 ***
Typass       0.0932763  0.0791100   1.179   0.2414    
age         -0.0005335  0.0012852  -0.415   0.6790    
Sexe         0.0951009  0.0499814   1.903   0.0602 .  
dest        -0.1043560  0.0692713  -1.506   0.1353    
---
Signif. codes:  0 '***' 0.001 '**' 0.01 '*' 0.05 '.' 0.1 ' ' 1

(Dispersion parameter for Gamma family taken to be 0.04961642)

    Null deviance: 47.9536  on 99  degrees of freedom
Residual deviance:  5.0718  on 93  degrees of freedom
AIC: 1817.9

Number of Fisher Scoring iterations: 5
\end{verbatim}
}
\end{frame}

\begin{frame}{}

\begin{center}
\includegraphics[width = 6.5cm]{figures/MYdiag.pdf}
\end{center}

Valeurs aberrantes ! (voir chapitre~\ref{robust.section})
\end{frame}

\subsection{Exemple 5 : Étude CARDIA}

\begin{frame}{Exemple 5 : Étude CARDIA}

Données de l'étude CARDIA (Coronary Artery Risk Development in young
Adults), disponibles depuis les archives de données du JASA à
\url{http://lib.stat.cmu.edu/jasadata/}, voir \citeN{Prei:Gale:Lohm:Wage:2000}.

\label{CARDIApage}
\bigskip

5\,115 adultes âgés de 18 à 30 ans ont été suivis et examinés jusqu'à 5 fois de
1986 à 1996, aux années 0, 2, 5, 7 et 10.

\bigskip

On considère le sous-échantillon de 3693 individus pour lesquels l'information complète est disponible
pour les 4 premières visites (aucune valeur manquante).

\bigskip

Variable réponse (à chaque visite) : statut tabagique autodéclaré (oui/non, codé 1/0).
\end{frame}


\begin{frame}{}

Prédicteurs (à chaque visite) :
\begin{itemize}
\item âge (dépendant du temps)
\item cohorte de naissance (1955--1962, 1963--1967) (indépendant du temps)
\item éducation (secondaire ou moins, études collégiales partielles, diplôme universitaire) (dépendant/(in)dépendant du temps)
\item race-sexe (hommes noirs, femmes noires, hommes blancs, femmes blanches) (indépendant du temps)
\end{itemize}

\vfill
{\color{red} $\bigstar$ Problèmes de données manquantes (missingness), voir \citeN{Prei:Gale:Lohm:Wage:2000}}.
\end{frame}

\begin{frame}[fragile]{}
{\tiny
\begin{verbatim}
Call:
gee(formula = smoke ~ age + factor(birth) + factor(education) + 
    factor(racesex), id = id, data = CARDIA.sub, family = binomial, 
    corstr = "exchangeable")

Summary of Residuals:
       Min         1Q     Median         3Q        Max 
-0.5506120 -0.3121130 -0.1189440  0.4816228  0.9095992 

Coefficients:
                      Estimate Naive S.E.    Naive z Robust S.E.   Robust z
(Intercept)        -1.33707777 0.54231807  -2.465486  0.54510289  -2.452891
age                 0.06480087 0.02646754   2.448315  0.02669744   2.427231
factor(birth)2     -0.19856466 0.14606724  -1.359406  0.14563751  -1.363417
factor(birth)3     -0.40380459 0.23689643  -1.704562  0.23858001  -1.692533
factor(education)2 -0.68712005 0.08313860  -8.264754  0.08490269  -8.093030
factor(education)3 -1.96140806 0.10016299 -19.582164  0.10054908 -19.506972
factor(racesex)2   -0.17668132 0.09763816  -1.809552  0.09900202  -1.784623
factor(racesex)3   -0.13445604 0.10634023  -1.264395  0.10833799  -1.241079
factor(racesex)4   -0.12730906 0.10530815  -1.208919  0.10599195  -1.201120

Estimated Scale Parameter:  1.000362
Number of Iterations:  1

Working Correlation
         [,1]     [,2]     [,3]     [,4]
[1,] 1.000000 0.727139 0.727139 0.727139
[2,] 0.727139 1.000000 0.727139 0.727139
[3,] 0.727139 0.727139 1.000000 0.727139
[4,] 0.727139 0.727139 0.727139 1.000000
\end{verbatim}
}
\end{frame}


\subsection{Exemple 6 : Abondance d'oiseaux}

\begin{frame}{Exemple 6 : Abondance d'oiseaux}

Les dénombrements d'oiseaux sont-ils en déclin? Abondance d'espèces animales.

\bigskip

L'exemple de 5 espèces de Seal Island (Nouvelle-Écosse,
Canada) qui hivernent principalement dans le sud des États-Unis : Northern Flicker, Red-wing Blackbird, Rusty Blackbird,
Ruby-crowned Kinglet, Yellow-rumped Warbler.

\bigskip

Données fournies par le prof.\ I. MacLaren, Dalhousie University.

\bigskip

2245 observations (saison automnale, entre le 1\textsuperscript{er} août et le 14 novembre, sur 40 ans
depuis 1963), soit 449 jours d'observations pour chaque espèce d'oiseau.

\end{frame}

\begin{frame}{}
\begin{tabular}{ccc}
\includegraphics[width = 3.5cm]{figures/novascotia.jpg}   & \includegraphics[width = 2.5cm]{figures/northernflicker.jpg}  & \includegraphics[width = 2.5cm]{figures/redwing2.jpg} \\
 \includegraphics[width = 2.5cm]{figures/rustyblackbird.jpg}  & \includegraphics[width = 2.5cm]{figures/ruby_crowned_kinglet.jpg} & \includegraphics[width = 2.5cm]{figures/Yellow_rumped_Warbler.jpg}
\end{tabular}

\end{frame}

\begin{frame}{}

Variable réponse : nombre d'oiseaux observés (décomptes, \texttt{NUMBER}).

\bigskip
Prédicteurs :
\begin{itemize}
\item année (\texttt{YEAR}) -- {\color{blue}Intérêt principal}.
\item jour (\texttt{DAY})
\item nombre
d'observateurs sur l'île (\texttt{NO.OBS})
\item direction du vent à minuit la veille (\texttt{WIND.0})
\item direction du vent à midi la veille (\texttt{WIND.1})
\item direction du vent à minuit deux jours avant (\texttt{WIND.2})
\item direction du vent à midi deux jours avant (\texttt{WIND.3})
\item vitesse du vent ce jour-là (\texttt{WINDSPD})
\item condition du ciel (1=dégagé, 2=nuageux, 3=brouillard, 4=pluie) (\texttt{SKY})
\item proportion de lune illuminée (\texttt{MOON})
\item espèce d'oiseau (1--5) (\texttt{BIRD})
\end{itemize}

\end{frame}


\begin{frame}[fragile]{}
Modèle de Poisson avec lien logarithmique :
\begin{eqnarray*}
\lefteqn{\texttt{NUMBER} \sim  \alpha + f_1(\texttt{YEAR}) + f_2(\texttt{DAY}) +
f_3(\texttt{NO.OBS})} \\
& + & f_4(\texttt{WIND.0}) + f_5(\texttt{WIND.1}) +
f_6(\texttt{WIND.2}) \\
& + & f_7(\texttt{WIND.3}) + f_8(\texttt{WINDSPD}) +
\mbox{factor(\texttt{SKY})}\\
& + & f_9(\texttt{MOON}) + \mbox{factor}(\texttt{BIRD})
\end{eqnarray*}

\bigskip

{\scriptsize
\begin{verbatim}
> require(mgcv)
> GrA.fit <- gam(NUMBER~s(YEAR,k=20)+s(DAY,k=20)+s(NO.OBS,k=8)+
   s(WIND.0,k=20)+s(WIND.1,k=20)+s(WIND.2,k=20)+s(WIND.3,k=20)+
   s(WINDSPD,k=20)+factor(SKY)+s(MOON,k=20)+factor(BIRD), 
   family=poisson,data=GrA,na.action=na.omit)
> plot(GrA.fit,scale=0,pages=1,pers=T,all.terms=T,cex.lab=2)
\end{verbatim}}
\end{frame}

\begin{frame}{}

\begin{center}
\includegraphics[width = 6.5cm]{figures/GrA.pdf}
\end{center}
\vfill{\color{red} $\bigstar$ Excès de zéros possible. Revue : \citeN{Rido:Deme:Hind:1998}}
\end{frame}

\subsection{Résumé}

\begin{frame}{Qu'ont en commun les exemples 1 à 6?}

{\color{blue}
Même formulation générale du modèle :
$$g(E(Y_i)) = g(\mu_i) = $$
$$= \eta_i = \left\{
   \begin{array}{l}
   \beta_0 + \beta_1 x_{i1}+\ldots+\beta_p x_{ip} \\
  \alpha + f_1(x_{i1})+ \ldots + f_p(x_{ip})
  \end{array} \right.$$}

\bigskip

$g$ est appelée la fonction de lien (link function).

\bigskip

$\eta_i$ est appelé le prédicteur linéaire (additif).

\end{frame}

\begin{frame}{Notation}

{\color{blue}Variable réponse} : \\
observations $y_1, \ldots, y_n$
pour $n$ individus (ou sujets) correspondant aux variables aléatoires $Y_1,
\ldots, Y_n$.

\bigskip

$y_i$ (resp.\ $Y_i$) peut être multidimensionnel, auquel cas
on notera ses éléments $y_{it}$ avec $t=1,\ldots,n_i$.

\bigskip

{\color{blue}Propriétés distributionnelles de $Y_i$} :

$Y_i$ issu d'une distribution $F_{\theta_i}$ :
$$Y_i \sim F_{\theta_i}$$
On note l'espérance $E(Y_i) = \mu_i$ et la variance $Var(Y_i) = v_i$.
\end{frame}

\begin{frame}{}

{\color{blue}Variables explicatives (prédicteurs)} :
\begin{itemize}
\item pour $y_i$ unidimensionnel :
$$x_i = (x_{i0}, x_{i1}, \ldots, x_{ip})^T$$
est un vecteur de dimension $(p+1)$. $x_{i0}$ est habituellement égal à 1
(ordonnée à l'origine). Information combinée :
$$X =  \left( \begin{array}{c}
x_1^T\\
\vdots \\
x_n^T
\end{array} \right) =
\left(
\begin{array}{cccc}
x_{10} & x_{11} & \cdots & x_{1p}\\
x_{20} & x_{21} & \cdots & x_{2p}\\
\vdots & \vdots & \vdots & \vdots \\
x_{n0} & x_{n1} & \cdots & x_{np}\\
\end{array} \right)
$$

\label{Ximatrix}
\end{itemize}
\end{frame}

\begin{frame}{}

\begin{itemize}
\item pour $y_i$ multidimensionnel :
$$X_i = \left(
\begin{array}{cccc}
x_{i10} & x_{i11} & \cdots & x_{i1p}\\
x_{i20} & x_{i21} & \cdots & x_{i2p}\\
\vdots & \vdots & \vdots & \vdots \\
x_{in_i0} & x_{in_i1} & \cdots & x_{in_ip}\\
\end{array} \right)$$
\end{itemize}

Information combinée :
$$X = \left(
\begin{array}{c}
X_1 \\
X_2 \\
\vdots\\
X_n
\end{array} \right) $$
\end{frame}

\begin{frame}{Connaissances préalables}

Familiarité avec les concepts suivants :
\begin{itemize}
\item Estimation (maximum de vraisemblance, moindres carrés, méthode des moments, etc.)
\item Intervalles de confiance.
\item Distributions d'échantillonnage.
\item Tests d'hypothèses (hypothèse nulle et alternative, statistique de test et
sa distribution, valeur-$p$, seuil de signification, etc.).
\item Régression linéaire.
\item Calcul différentiel et notation matricielle.
\item $\ldots$
\end{itemize}

\end{frame}


\begin{frame}[allowframebreaks]{Principes de la modélisation statistique}

\bigskip

Une analyse statistique devrait comprendre les étapes suivantes :
\begin{enumerate}
\item {\color{blue}Analyse exploratoire des données} pour vérifier la qualité des données et aider
à la formulation du modèle. Considérer chaque variable séparément et se poser les
questions suivantes :
\begin{itemize}
\item Quelle est l'échelle de mesure?
\item Quelle est la forme de la distribution?
\item Comment est-elle associée aux autres variables?
\end{itemize}
\item {\color{blue}Formulation du modèle} :
  \begin{itemize}
  \item Distribution de probabilité de $Y$.
  \item Lien avec les variables explicatives.
  \end{itemize}
\item {\color{blue}Estimation des paramètres} : quelle méthode? (maximum de vraisemblance,
moindres carrés, méthode des moments, etc.)
\newpage
\item {\color{blue}Résidus et vérification du modèle} :
  \begin{itemize}
  \item Inspection graphique des résidus (indépendance, normalité approximative, moyenne nulle, variance constante).
  \item Tests d'adéquation (goodness of fit).
  \end{itemize}
\item {\color{blue}Inférence et interprétation} :\\
Tester les hypothèses sur les paramètres pour obtenir un modèle parcimonieux. Puis (et seulement
alors) interpréter les paramètres.
\end{enumerate}

\end{frame}

\begin{frame}{Remerciements}
Ces diapositives ont été inspirées par le cours GLAM de Eva Cantoni de l'Université de Genève.
\end{frame}

\begin{frame}[allowframebreaks]{Références du chapitre}
\tiny
\printbibliography[heading=none,section=\therefsection]
\end{frame}

\end{refsection}

\end{document}
